\documentclass[11pt,a4paper]{article}
\usepackage[utf8]{inputenc}
\usepackage[sfdefault]{noto}
\usepackage[T1]{fontenc}
\usepackage[spanish]{babel}
\usepackage[left=2.5cm,right=2.5cm,top=3cm,bottom=3cm]{geometry}
\usepackage{hyperref}
\usepackage{tabularx}
\usepackage{booktabs}
\usepackage{amsmath}
\usepackage{enumitem}
\usepackage{xcolor}
\usepackage{titlesec}

% Colores y estilos
\definecolor{ucmblue}{RGB}{0, 51, 102}
\titleformat{\section}{\color{ucmblue}\Large\bfseries}{\thesection}{1em}{}
\titleformat{\subsection}{\color{ucmblue}\large\bfseries}{\thesubsection}{1em}{}

\title{\vspace{-2cm}\color{ucmblue}\textbf{Tarea Práctica: Predicción de Tiempos de Sentencia en el Poder Judicial Chileno}}
\author{Sergio Hernández\\MDS211 - Machine Learning\\Universidad Católica del Maule}
\date{}

\begin{document}

\maketitle

\section{Objetivo de la Tarea}
El objetivo de esta evaluación es desarrollar, optimizar y comparar dos enfoques de Machine Learning (Redes Neuronales y Gradient Boosting) para predecir la duración de los procesos judiciales (tiempos de sentencia) utilizando datos reales del Poder Judicial de Chile. Además de la predicción puntual, se requiere una \textbf{evaluación rigurosa de la incertidumbre} asociada a dichas predicciones.

\section{Descripción del Problema y Datos}
Se utilizarán las estadísticas oficiales del Poder Judicial \footnote{\url{https://numeros.pjud.cl/Estadisticas/Sentencias}}. Los estudiantes deberán recolectar (o descargar) un subconjunto representativo de causas terminadas, utilizando atributos como tipo de tribunal, competencia, materia, región, entre otros.

La variable objetivo ($y$) será el \textbf{tiempo de sentencia} (por ejemplo, calculado como la diferencia en días entre la fecha de ingreso y la fecha de término o sentencia).

\section{Requerimientos y Etapas del Proyecto}

\subsection{1. Recopilación y Preprocesamiento de Datos}
\begin{itemize}
    \item \textbf{Limpieza:} Tratamiento de valores nulos, registros duplicados y \textit{outliers} en los tiempos de sentencia (ej. causas anómalamente largas o con errores de digitación).
    \item \textbf{Ingeniería de Características:} Codificación de variables categóricas (ej. \textit{One-Hot Encoding} o \textit{Feature Hashing} para competencias y materias), transformación de variables temporales y escalado/normalización para el entrenamiento de la Red Neuronal.
    \item \textbf{Partición:} Separación de datos en conjuntos de Entrenamiento, Validación y Prueba (se sugiere considerar una partición temporal si aplica).
\end{itemize}

\subsection{2. Modelado y Ajuste de Hiperparámetros}
Deberá entrenar y optimizar los siguientes modelos:
\begin{itemize}
    \item \textbf{Modelo 1: Gradient Boosting.} Puede utilizar implementaciones como XGBoost, LightGBM o CatBoost. Realice una búsqueda de hiperparámetros (ej. tasa de aprendizaje, profundidad del árbol, número de estimadores) utilizando validación cruzada.
    \item \textbf{Modelo 2: Red Neuronal (Deep Learning).} Diseñe un perceptrón multicapa (MLP) utilizando frameworks como JAX/Flax o PyTorch. Ajuste hiperparámetros como el número de capas, neuronas, funciones de activación, \textit{learning rate} y \textit{batch size}.
\end{itemize}

\subsection{3. Evaluación de la Incertidumbre}
Dado que el tiempo de resolución de una causa judicial es altamente estocástico, la predicción puntual es insuficiente.
\begin{itemize}
    \item Para \textbf{Gradient Boosting}: Implemente regresión por cuantiles (\textit{Quantile Regression}) para obtener intervalos de predicción (ej. cuantiles 10\% y 90\%).
    \item Para la \textbf{Red Neuronal}: Implemente una red MLP para estimar los parámetros de una distribución (ej. \textit{Gaussian Log-Likelihood} loss) o utilice librerías de inferencia probabilística como NumPyro \footnote{\url{https://num.pyro.ai/en/stable/getting_started.html}}.
\end{itemize}

\subsection{4. Análisis y Conclusiones}
\begin{itemize}
    \item Compare las métricas de error puntual (RMSE, MAE).
    \item Compare la calidad de los intervalos de incertidumbre (ej. Cobertura empírica vs. Cobertura nominal, Amplitud promedio del intervalo).
    \item Discuta qué modelo ofrece una mejor calibración y cuál recomendaría para su implementación en un sistema de apoyo a la gestión de tribunales.
\end{itemize}

\newpage
\section{Rúbrica de Evaluación}

La tarea será evaluada con base en la siguiente rúbrica (Puntaje Total: 100 puntos):

\vspace{0.5cm}
\noindent
\begin{tabularx}{\textwidth}{>{\hsize=0.2\hsize}X >{\hsize=0.25\hsize}X >{\hsize=0.25\hsize}X >{\hsize=0.3\hsize}X}
\toprule
\textbf{Criterio} & \textbf{Insuficiente (0-5 pts)} & \textbf{Adecuado (6-12 pts)} & \textbf{Excelente (13-20 pts)} \\
\midrule
\textbf{1. Preprocesamiento} & 
Manejo deficiente de nulos, sin transformación de variables ni tratamiento de outliers. & 
Limpieza básica, codificación estándar, pero escalado o ingeniería de características incompleta. & 
Preprocesamiento exhaustivo, justificación de exclusión de outliers y excelente ingeniería de características. \\
\addlinespace
\textbf{2. Implementación de Modelos} & 
Modelos mal configurados, \textit{data leakage} o errores graves en la compilación de la Red Neuronal. & 
Ambos modelos corren correctamente, pero la arquitectura de la red es muy simple o el GB no utiliza validación cruzada. & 
Implementación robusta (idealmente JAX/Flax para NN), uso de buenas prácticas y código optimizado. \\
\addlinespace
\textbf{3. Ajuste de Hiperparámetros} & 
Se utilizan hiperparámetros por defecto sin ninguna estrategia de búsqueda. & 
Búsqueda empírica básica (\textit{Grid Search} pequeño) sin métricas claras de validación. & 
Búsqueda sistemática (ej. \textit{Random Search} o \textit{Grid Search}) evaluada con validación cruzada. \\
\addlinespace
\textbf{4. Evaluación de Incertidumbre} & 
No se evalúa la incertidumbre, solo se reportan predicciones puntuales. & 
Implementación teórica básica de cuantiles o MC Dropout, pero intervalos mal interpretados o no calculados correctamente. & 
Excelente formulación bayesiana o cuantílica. Cobertura de intervalos calculada y visualizada de forma precisa. \\
\addlinespace
\textbf{5. Análisis y Conclusiones} & 
Conclusiones superficiales, no se comparan adecuadamente los intervalos ni las métricas de error. & 
Comparación basada solo en MAE/RMSE. Breve mención a la incertidumbre sin profundidad analítica. & 
Análisis crítico profundo. Discusión clara sobre el equilibrio entre precisión puntual, calidad de la incertidumbre y costo computacional. \\
\bottomrule
\end{tabularx}

\vspace{1cm}
\noindent \textbf{Formato de Entrega:} Se debe entregar un documento en PDF con el informe metodológico y un repositorio (o archivo \texttt{.ipynb}) con el código documentado.

\end{document}
